% unidades/10-tame-ni-you-ni.tex
\chapter{\emoji{dart} \ruby{ために}{} vs \ruby{ように}{} — Propósito}
\unidadheader{10}{目的}{ために vs ように}{Expresar metas concretas y estados deseados}

\begin{tcolorbox}[vocabbox, title={\emoji{pushpin} Vocabulario}]
\begin{tabularx}{\linewidth}{llX}
\toprule
\textbf{Japonés} & \textbf{Español} & \textbf{Tipo} \\ \midrule
\vocabitem{合格する}{ごうかくする}{aprobar (examen)}{v. suru}
\vocabitem{痩せる}{やせる}{adelgazar}{v. G2}
\vocabitem{貯金する}{ちょきんする}{ahorrar dinero}{v. suru}
\vocabitem{健康になる}{けんこうになる}{ponerse sano}{v. compuesto}
\vocabitem{努力する}{どりょくする}{esforzarse}{v. suru}
\vocabitem{覚える}{おぼえる}{memorizar / aprender}{v. G2}
\vocabitem{泳ぐ}{およぐ}{nadar}{v. G1}
\bottomrule
\end{tabularx}
\end{tcolorbox}

\subsection*{\emoji{speech-balloon} Frases Clave}
\frase{\ruby{医者}{いしゃ}になるために\ruby{毎日}{まいにち}\ruby{勉強}{べんきょう}しています。}{Estudio todos los días para llegar a ser médico.}
\frase{\ruby{泳}{およ}げるように\ruby{毎日}{まいにち}\ruby{練習}{れんしゅう}しています。}{Practico todos los días para poder nadar.}
\frase{\ruby{忘}{わす}れないようにメモを\ruby{取}{と}ります。}{Tomo notas para no olvidar.}
\frase{\ruby{健康}{けんこう}のために\ruby{毎朝}{まいあさ}\ruby{走}{はし}っています。}{Corro cada mañana para mi salud.}

\subsection*{\emoji{gear} Gramática}
\patron{ために — meta concreta directa}
  {[V-dict / N+の] + ために}
  {\ej{\ruby{大学}{だいがく}に\ruby{入}{はい}るために\ruby{勉強}{べんきょう}した。}{Estudié para entrar a la universidad.}
   \ej{\ruby{家族}{かぞく}のために\ruby{働}{はたら}いています。}{Trabajo para mi familia.}
   \ej{\ruby{日本語}{にほんご}を\ruby{覚}{おぼ}えるために\ruby{毎日}{まいにち}\ruby{練習}{れんしゅう}します。}{Practico todos los días para aprender japonés.}}

\patron{ように — capacidad / estado deseado}
  {[V-potential / V-negative] + ように}
  {\ej{\ruby{泳}{およ}げるように\ruby{練習}{れんしゅう}している。}{Practico para poder nadar (meta: capacidad)}
   \ej{\ruby{忘}{わす}れないようにメモした。}{Tomé notas para no olvidar (meta: estado negativo evitado)}
   \ej{\ruby{風邪}{かぜ}を\ruby{引}{ひ}かないように\ruby{気}{き}をつけて。}{Cuídate para no resfriarte.}}

\comparacion{ために vs ように}
  {ために\\\small meta concreta e intencional\\\small \ruby{医者}{いしゃ}になるために\\\textit{Para convertirme en médico}}
  {ように\\\small capacidad o estado deseado\\\small \ruby{泳}{およ}げるように\\\textit{Para poder nadar}}
  {ために se usa con verbos de acción volicional en diccionario. ように se usa con verbos potenciales, negativos, o cuando el sujeto principal no controla el resultado. Nunca se dice 「医者になれるように」con ため si hay control directo.}

\coloquial{〜ために / 〜ように}{natural en cualquier registro}{
  Ambas formas son naturales en habla cotidiana. En casual: 〜ようにする (procuro hacer) es muy frecuente: 毎日運動するようにしている。}

\culturabox{\ruby{努力}{どりょく}の\ruby{文化}{ぶんか} — La cultura del esfuerzo}{
  En Japón, el esfuerzo sostenido (\ruby{努力}{どりょく}) y la perseverancia (\ruby{忍耐}{にんたい}) son virtudes culturales fundamentales. ために y ように reflejan esta orientación hacia la meta: el japonés siempre quiere saber «¿para qué?» en una acción. \ruby{目標}{もくひょう}のために\ruby{頑張}{がんば}る es casi un mantra.
}

\lectura{\ruby{夢}{ゆめ}のために\\[0.4em]
  アナは\ruby{日本語}{にほんご}の\ruby{先生}{せんせい}になるために、\ruby{毎日}{まいにち}\ruby{勉強}{べんきょう}している。\ruby{漢字}{かんじ}を\ruby{忘}{わす}れないようにカードを\ruby{作}{つく}り、\ruby{話}{はな}せるようになるために\ruby{日本人}{にほんじん}の\ruby{友達}{ともだち}と\ruby{会話}{かいわ}の\ruby{練習}{れんしゅう}をしている。\ruby{夢}{ゆめ}のために\ruby{努力}{どりょく}することが\ruby{大切}{たいせつ}だと\ruby{信}{しん}じている。}
{Por el sueño\\[0.4em]
  Ana estudia todos los días para convertirse en profesora de japonés. Hace tarjetas para no olvidar los kanji, y practica conversación con amigos japoneses para poder hablar. Cree que es importante esforzarse por el propio sueño.}

\begin{tcolorbox}[ejerciciobox, title={\emoji{pencil} Ejercicios}]
\begin{enumerate}
  \item ¿ために o ように? \ruby{日本語}{にほんご}が\ruby{上手}{じょうず}になる\_\_\_\ruby{勉強}{べんきょう}します。
  \item ¿ために o ように? \ruby{日本語}{にほんご}が\ruby{話}{はな}せる\_\_\_\ruby{練習}{れんしゅう}します。
  \item Crea una oración con ように describiendo una meta personal.
\end{enumerate}
\textbf{Respuestas:}
1) ために (meta con verb de cambio de estado volicional) \quad
2) ように (capacidad: 話せる = potencial) \quad
3) (libre)
\end{tcolorbox}

\begin{tcolorbox}[kanjibox, title={\emoji{mahjong} Kanjis}]
\begin{tabularx}{\linewidth}{cllXX}
\toprule
\textbf{Kanji} & \textbf{On} & \textbf{Kun} & \textbf{Significado} & \textbf{Vocab} \\ \midrule
\kanjirow{合}{ごう}{\ruby{あ}{う}}{coincidir/pasar}{\ruby{合格}{ごうかく} / \ruby{合}{あ}う}
\kanjirow{格}{かく}{—}{rango/calidad}{\ruby{合格}{ごうかく} / \ruby{格差}{かくさ}}
\kanjirow{努}{ど}{\ruby{つと}{める}}{esforzarse}{\ruby{努力}{どりょく} / \ruby{努}{つと}める}
\kanjirow{力}{りょく・りき}{\ruby{ちから}{}}{fuerza}{\ruby{努力}{どりょく} / \ruby{力}{ちから}}
\bottomrule
\end{tabularx}
\end{tcolorbox}
